\section{$R^2$ 衡量了什么、不衡量什么？}\label{sec:r-squared}

\subsection{平方和分解}\label{sec:sst-decomp}

在含截距的样本内 OLS 中，\eqref{eq:ols-slope} 所在推导的一阶条件给出 $\sum_i\hat u_i=0$ 和 $\sum_ix_i\hat u_i=0$，其中 $\hat u_i=y_i-\hat y_i$。
由于 $\hat y_i=\hat\beta_0+\hat\beta_1x_i$，有 $\sum_i(\hat y_i-\bar y)\hat u_i=0$。

将 $y_i-\bar y=(\hat y_i-\bar y)+\hat u_i$ 平方后求和，交叉项为零，得到

\begin{equation}
  \underbrace{\sum_i(y_i-\bar y)^2}_{SST}
  =\underbrace{\sum_i(\hat y_i-\bar y)^2}_{SSE}
  +\underbrace{\sum_i\hat u_i^2}_{SSR}.
\end{equation}

当 $SST>0$ 时，拟合优度定义为被解释平方和占总平方和的比例：

\begin{equation}
  R^2=\frac{SSE}{SST}=1-\frac{SSR}{SST}.\label{eq:r-squared}
\end{equation}

在上述条件下，$0\leq R^2\leq1$；残差平方和相对总平方和越小，样本内拟合程度越高。

\subsection{常见误读}\label{sec:r2-misread}

$R^2$ 只比较同一样本中的平方和，不检验解释变量是否与结构误差相关，也不识别因果方向，因此高 $R^2$ 不能单独证明设定正确或因果关系成立。

对同一组观测增加解释变量时，OLS 可以把新系数设为零而保留原拟合，故最小残差平方和不会增大，普通 $R^2$ 不会下降；这一性质本身不保证新增变量有因果意义。斜率估计的精度则应看标准误，不能仅由 $R^2$ 高低判断。
